\documentclass[10pt,a4paper]{amsart}
\usepackage[margin=19mm]{geometry}
\usepackage{amsmath,amssymb,mathtools}
\usepackage[scr=rsfs,cal=euler]{mathalfa}
\usepackage{xfrac}
\usepackage[unicode,hidelinks,bookmarks=false]{hyperref}
\newcommand{\ie}{i.e.\ }
\newcommand{\cf}{cf.\ }
\newcommand{\resp}{respectively\ }
\newcommand{\Subsection}{\subsection}
\providecommand{\nobreakdash}{\nobreak\hbox{-}}
\setlength{\emergencystretch}{2em}
\multlinegap0pt
\usepackage{amssymb}
\usepackage[shortlabels]{enumitem}
\usepackage{dsfont}
\newtheorem{theorem}{Theorem}
\newtheorem{lemma}{Lemma}
\newcommand{\mutau}{\mu_{\tau}}
\title{Exponential Bounds and Analyticity for the Tree Builder Random Walk}
\author{Caio Alves, Rodrigo Ribeiro}
\date{Source: arXiv:2603.28578v1 (CC BY 4.0)}
\begin{document}
\maketitle

\noindent\emph{Source and licence.} Exponential Bounds and Analyticity for the Tree Builder Random Walk. Version arXiv:2603.28578v1 (CC BY 4.0), distributed by arXiv under the Creative Commons Attribution 4.0 International licence, http://creativecommons.org/licenses/by/4.0/, as recorded in the arXiv licence metadata for this exact version and shown on the abstract page https://arxiv.org/abs/2603.28578v1. Full licence notice: \texttt{licenses/arxiv-2603.28578-CC-BY-4.0.txt}.\par\smallskip

\noindent\textit{Editorial scope.} Counted unit: the article's lemma l:N\_ld\_bound (source0.tex lines 533--546). The article's complete original proof of it is delivered with it (source0.tex lines 547--602, one \texttt{proof} environment of 56 lines). No mathematics is written, repaired or re-derived: the statement and the proof are the article's own lines, in the article's own notation, and no retained line is substituted or re-flowed.  Retained uncounted as the source-local context the proof needs: lines 125--130; lines 502--516.  Nothing was omitted from any retained span, so the package carries no omission record.  No retained line contains a citation, so the wrapper carries no bibliography and no bibliography entry.  No retained line contains a figure environment, an image-inclusion command or drawing code, so no image is delivered and the package declares no asset dependency.  Editorial formatting only, in addition: an AMS \texttt{amsart} preamble that restates the article's own declarations and macros used by the retained lines, namely the environments \texttt{theorem}, \texttt{lemma} on separate counters, together with the standard packages those lines need.  The retained environments print in this document as Theorem 1, Lemma 1 in order of appearance, whereas the article prints the same items with the numbering of its own full document.  The complete native source of the article as distributed in the arXiv e-print for this exact version is bundled unchanged as \texttt{sources/197457/source0.tex}.
\begin{theorem}[Exponential tail of $\tau_1$]\label{t:tautail} Consider a TBRW with leaf distribution $Q$ and fix $\kappa\in (0,1]$, then there exists a constant $c$ depending on $\kappa$ only, such that
$$
\sup_{Q \in \mathcal{Q}_\kappa}P_Q(\tau_1 \ge n) \le e^{-cn}.
$$
    
\end{theorem}

\begin{lemma}
\label{l:tau_ld_bound}
Fix $\kappa \in (0,1]$. Then, for all $\varepsilon > 0$, there exists $c = c(\varepsilon, \kappa)$ such that for all $m \in \mathbb{N}$ the following bound holds
    \begin{equation}
       \sup_{Q\in \mathcal{Q}_\kappa} P_Q \left(
                \left|
                    \tau_m - m \mutau
                \right|
                    >
                \varepsilon m
            \right)
            \leq
        e^{-c(\varepsilon) m}.
    \end{equation}  
\end{lemma}

\begin{lemma}
    \label{l:N_ld_bound} Fix $\kappa \in (0,1]$. Then, for all $\varepsilon >0$, there exists a positive constant $c = c(\varepsilon,\kappa)$ such that
    \begin{equation}
       \sup_{Q\in \mathcal{Q}_\kappa} P_Q \left(
                \left|
                    N(t) - \frac{t}{\mutau}
                \right|
                    >
                \varepsilon t
            \right)
            \leq
        e^{- c t}.
    \end{equation}  
\end{lemma}

\begin{proof}
    We start by defining
    \begin{equation}
        \label{eq:mplusdef}
        m_+ = m_+(t, \varepsilon)
            = 
            \left(
                \frac{1}{\mutau} + 
                \varepsilon
            \right)
            t
    \end{equation}
    and therefore,
    \begin{equation}
        \label{eq:Ntldimplication}
        N(t) 
            >
        \left(
            \frac{1}{\mutau} + 
            \varepsilon
        \right)
        t
            \implies
                \tau_{\lceil m_+ \rceil} \leq t 
               = \mutau \lceil m_+ \rceil\left(1-\varepsilon\mutau + O(\varepsilon^2\mutau^2) \right),
    \end{equation}
    which is possible since, by Theorem \ref{t:tautail} and the definition of $\tau_1$, for all $Q \in \mathcal{Q}_\kappa$, $\mutau$ is bounded from above and from below by positive constants depending on $\kappa$ only. Finally, using Lemma \ref{l:tau_ld_bound}, we can make $\varepsilon$ small enough so that
    \begin{equation}
        \label{eq:Ntldbound2}
        P_Q \left(
            N(t) 
                >
            \left(
                \frac{1}{\mutau} + 
                \varepsilon
            \right)
            t
        \right)
        \leq
            \exp
            \left\{
                -c(\varepsilon) m_+
            \right\}
        =
            \exp
            \left\{
                -c(\varepsilon)  
                \left(
                    \frac{1}{\mutau} + 
                    \varepsilon
                \right)
                t
            \right\},
    \end{equation}
    since the events $\{N(t) > m_+\}$ and $\{\tau_{\lceil m_+ \rceil} \leq t\}$ are the same. The proof for the upper bound on $N(t)$ follows in the same manner.
\end{proof}

\end{document}
